\par\bigskip
\begin{center}\pdfbookmark[2]{7. Seculare und periodische Störungen}{ceres-007}[7.]\end{center}
\noindent\hspace*{\parindent}Aus obigen Formeln folgt nunmehr

I) Secularer Zuwachs der Excentricität \(=\tfrac12e'a\mathfrak C^{(0)}nt\sin(\tilde\omega'-\tilde\omega)\); [nun ist] \(\mathfrak C^{(0)}=\tfrac12Q''\), also der Zuwachs \(=\tfrac14e'aQ''nt\sin(\tilde\omega'-\tilde\omega)\)
\[
\begin{array}{cc}
\text{Für Ceres ist:}&\text{[Für Jupiter:]}\\[.5ex]
\tilde\omega=352^{0}\,58'\,0''&\tilde\omega'=191^{0}\,10'\,20''\\[.5ex]
\log n\ \text{[täglich]}=2{,}886\,2959&\log e'=8{,}681\,3317\,[-10]
\end{array}
\]
[und hieraus findet man
\[
\text{jährliche Abnahme der Excentricität]}\quad0{,}00000\,59089.
\]

II) Secularzunahme der Länge der Sonnenferne:
\[
=\frac a4\left(B^{(0)}-2\mathfrak B^{(0)}-2\frac{e'}e\mathfrak C^{(0)}\cos(\tilde\omega'-\tilde\omega)\right)nt.
\]
Nun ist
\[
B^{(0)}-2\mathfrak B^{(0)}=Q',
\]
also jene Zunahme
\[
=\frac a4\left(Q'-\frac{e'}eQ''\cos(\tilde\omega'-\tilde\omega)\right)nt.
\]

Für Ceres, Motus Annuus Apsidum \(=+70''{,}1503.\)

Nach diesen Secularänderungen der Excentricität und Sonnenferne fallen die oben roth\refstepcounter{footnote}\hyperlink{ceres7-note}{[*)]}\begingroup
\makeatletter
\long\def\@makefntext#1{\noindent[*) #1}
\makeatother
\footnotetext{\hypertarget{ceres7-note}{}Seite 389, wo sie im Abdruck ebenfalls unterstrichen worden sind.]}
\endgroup{} unterstrichenen Gleichungen ganz weg.

III) Der beständige Theil von \(\delta r\) ist \(=[-aa(2g+\mathfrak A^{(0)})=]+\tfrac16aaB^{(0)}\)
\[
\text{[also]}\quad=-0{,}000\,9475.
\]

IV) Der Coefficient von \(\cos iD\) [in \(\delta r\)] wird (weil \(i\) oben sowohl die positiven als die negativen Zahlen bedeutete und nun auf die positiven eingeschränkt wird)
\[
\left[\frac{2aa\mathfrak A^{(i)}}{hh-1}=\right]\frac{2aa}{hh-1}\left(\frac ihA^{(i)}+\tfrac12B^{(i)}\right)\mathpunct{\cdot}
\]

\hypertarget{native-398}{}[In unserm Falle, \(\log n'=2{,}475\,8562\)] wird \(\log\theta=9{,}589\,5603\) [und hiemit der Theil von \(\delta r\):]
\[
\begin{aligned}
&+0{,}001\,0304\cos D&&-0{,}000\,0350\cos5D\\
&-0{,}003\,8023\cos2D&&-0{,}000\,0128\cos6D\\
&-0{,}000\,4206\cos3D&&-0{,}000\,0050\cos7D.\\
&-0{,}000\,1077\cos4D
\end{aligned}
\]

V) In \(\delta v\) wird der Coefficient von \(\sin iD\)
\[
\begin{aligned}
&-\frac{4ai(1-\theta)}{ii(1-\theta)^2-1}\left(\frac1{1-\theta}A^{(i)}+\tfrac12B^{(i)}\right)\\
&+\frac a{i(1-\theta)}\left(\frac{3A^{(i)}}{1-\theta}+2B^{(i)}\right)
\end{aligned}
\]
[also in unserm Falle die entsprechenden Glieder:]
\[
\begin{aligned}
&-231''{,}94\sin D&&+3''{,}05\sin5D\\
&+496{,}68\sin2D&&+1{,}07\sin6D\\
&+44{,}15\sin3D&&+0{,}41\sin7D.\\
&+10{,}07\sin4D
\end{aligned}
\]

VI) Zum Behuf der von \(M\) und \(M'+D\) [abhängigen Glieder] dient folgende Untersuchung. [Bisher sind in \(\delta v\) die Integrationsconstanten, welche die Gleichung 11 enthält, unberücksichtigt geblieben; sie können beliebig gewählt werden, da sie sich mit den elliptischen Elementen \(e\) und \(\tilde\omega\) vereinigen, und sollen so bestimmt werden, dass die von \(M\) und \(D+M'=M+\tilde\omega-\bar{\tilde\omega}'\) abhängigen Glieder in \(\delta v\) verschwinden. Wenn die Integrationsconstanten \(f_1\) und \(f_1'\) so eingeführt werden, dass die Constante von \(\int rQ\sin V\mathord{.}\,n\,\mathrm{d}t\) gleich \(aaef_1+aae'f_1'\cos(\tilde\omega-\tilde\omega')\) und die aus \(\int rQ\cos V\mathord{.}\,n\,\mathrm{d}t\) entspringende gleich \(aae'f_1'\sin(\tilde\omega-\tilde\omega')\) wird, so ist der entsprechende Theil von \(\delta r\) mit Vernachlässigung der zweiten Potenzen der Excentricitäten
\[
\alpha)\qquad aaef_1\cos M+aae'f_1'\cos(M+\tilde\omega-\tilde\omega')
\]
und hiedurch entsteht in \(\delta v\) der Theil
\[
\beta)\qquad-2aef_1\sin M-2ae'f_1'\sin(M+\tilde\omega-\tilde\omega').]
\]

Die [gesamten von \(M\) und \(D+M'\) abhängigen] Glieder mit Inbegriff derer [\(\alpha\)], die aus den Constanten entspringen, seien
\hypertarget{native-399}{}
\[
\text{in }\delta r\ .\, .\, .\, .\quad aaef\cos M+aae'f'\cos(M+\tilde\omega-\tilde\omega')
\]
[so dass also
\[
\begin{aligned}
f&=f_1-\tfrac1{24}B^{(0)}-\tfrac14\mathfrak B^{(0)}\\
f'&=f_1'-\tfrac14\mathfrak C^{(0)}.
\end{aligned}
\]

Hiemit wird der gesamte von den beiden Argumenten abhängige Theil
\[
\begin{aligned}
\text{in }\delta v\ .\, .\, .\, .\quad&-2a\{f-\tfrac16B^{(0)}-\tfrac12\mathfrak B^{(0)}+\tfrac14\beta^{(0)}\}e\sin M\\
&-2a\{f'-\tfrac12\mathfrak C^{(0)}+\tfrac14\gamma^{(0)}\}e'\sin(M+\tilde\omega-\tilde\omega').]
\end{aligned}
\]
Man bestimmt \(f,f'\) so, dass der Inbegriff dieser Grössen in \(\delta v\) verschwindet. So bleibt
\[
\begin{aligned}
\text{in }\delta r\ .\, .\, .\, .\quad&aae(\tfrac12\mathfrak B^{(0)}-\tfrac14\beta^{(0)}+\tfrac16B^{(0)})\cos M\\
&+aae'(\tfrac12\mathfrak C^{(0)}-\tfrac14\gamma^{(0)})\cos(M+\tilde\omega-\tilde\omega')
\end{aligned}
\]
oder
\[
\begin{aligned}
&aae(\tfrac5{12}B^{(0)}+\tfrac14C^{(0)})\cos M\\
&+aae'(\tfrac14A'-\tfrac14C')\cos(M+\tilde\omega-\tilde\omega')
\end{aligned}
\]
[und numerisch]
\[
-0{,}000\,0526\cos M\qquad+0{,}000\,0141\cos(M+\tilde\omega-\tilde\omega').
\]

[Diese beiden Glieder lassen sich in eines vereinigen; bezeichnet man sie vorübergehend mit
\[
b_1\cos M+b_2\cos(M+\tilde\omega-\tilde\omega')
\]
und bestimmt man die Constanten \(c\) und \(C\) aus den Gleichungen
\[
\begin{aligned}
c\cos C&=b_1+b_2\cos(\tilde\omega-\tilde\omega')\\
c\sin C&=b_2\sin(\tilde\omega-\tilde\omega'),
\end{aligned}
\]
so erhalten sie die Form
\[
c\cos(M+C),
\]
numerisch] \(-0{,}000\,0633\cos(M+25^{0}\,56'\,46'').\)

VII) In \(\delta r\) wird der Coefficient von \(\cos(iD+M)\), [der mit] \(aaeE^{(i)}\) [bezeichnet werden mag]
\[
\begin{aligned}
aaeE^{(i)}&=aae\left\{\frac{\mathfrak B^{(i)}}{hh+2h}+\tfrac12\mathfrak A^{(i)}\left(\frac1{h+2}-\frac{4h}{hh-1}+\frac3h\right)\right\}\\
&=\frac{aae}{h(h+2)}\left\{\begin{aligned}
&\tfrac12C^{(i)}-\tfrac12B^{(i)}-\frac{3A^{(i)}}{1-\theta}\\
&-2\mathfrak A^{(i)}\left(\frac{3+ih(h-1)}{hh-1}\right)
\end{aligned}\right\}\mathpunct{\cdot}
\end{aligned}
\]

VIII) In \(\delta r\) wird der Coefficient von \(\cos((i+1)D+M')\), [der mit] \(aae'D^{(i)}\) \hypertarget{native-400}{}[bezeichnet werden mag]
\[
\begin{aligned}
aae'D^{(i)}&=\frac{aae'\mathfrak C^{(i)}}{hh+2h}\\
&=\frac{aae'}{hh+2h}\left\{\begin{aligned}
&\frac{i+1}{h+1}\bigl((2i+1)A^{(i+1)}-B^{(i+1)}\bigr)\\
&+(i+\tfrac12)B^{(i+1)}-\tfrac12C^{(i+1)}
\end{aligned}\right\}\mathpunct{\cdot}
\end{aligned}
\]

[Auch hier unterscheiden sich die Argumente \(iD+M\) und \((i+1)D+M'\) nur um constante Grössen; sie lassen sich auch in ein Glied vereinigen. Setzt man ähnlich, wie unter VI]
\[
\alpha\genfrac{}{}{0pt}{1}{\sin}{\cos}(x-A)+\beta\genfrac{}{}{0pt}{1}{\sin}{\cos}(x-B)
=f\genfrac{}{}{0pt}{1}{\sin}{\cos}(x-A+z),
\]
[so bestimmen sich \(f\) und \(z\) aus
\[
\begin{aligned}
\beta\sin(A-B)&=f\sin z\\
\alpha+\beta\cos(A-B)&=f\cos z.]
\end{aligned}
\]

IX) In \(\delta v\) wird der Coefficient von \(\sin(iD+M)\)
\[
\begin{aligned}
&=ae\left\{\begin{aligned}
&-2(h+1)E^{(i)}+\frac{2\mathfrak B^{(i)}-\tfrac12\beta^{(i)}}{h+1}\\
&-(2h+1)\frac{\mathfrak A^{(i)}}{hh-1}
\end{aligned}\right\}\\
&=\frac{ae}{h+1}\left\{\begin{aligned}
&-2E^{(i)}-2\mathfrak A^{(i)}\frac{\tfrac12i(h-1)-3}{hh-1}\\
&+\frac{(i+1)A^{(i)}}{1-\theta}
\end{aligned}\right\}\mathpunct{\cdot}
\end{aligned}
\]

X) In \(\delta v\) wird der Coefficient von \(\sin((i+1)D+M')\)
\[
\begin{aligned}
&=ae'\left\{-2(h+1)D^{(i)}+\frac{2\mathfrak C^{(i)}-\tfrac12\gamma^{(i)}}{h+1}\right\}\\
&=\frac{ae'}{h+1}\left\{\begin{aligned}
&-2D^{(i)}\\
&-\tfrac12\frac{i+1}{h+1}\bigl((2i+1)A^{(i+1)}-B^{(i+1)}\bigr)
\end{aligned}\right\}\mathpunct{\cdot}
\end{aligned}
\]
\par\medskip\centerline{\rule{.23\textwidth}{.35pt}}\medskip
